% =====================================================================
%  R. Nielsen, SAMFUNDETS NATURLOVE
%  In: For Ide og Virkelighed. Et Tidsskrift, udgivet af R. Nielsen, Bj. Bjørnson og Rud. Schmidt.  Aargang 1871. Første Bind (Bind 4).
%  Kjøbenhavn: Otto B. Wroblewskys Forlag; G. S. Wibes Bogtrykkeri.  Bind 4, pp. 289-315.  Antiqua.
%
%  TRANSCRIPTION -- GENERATED; DO NOT EDIT BY HAND.  Rasmus Nielsen's eight essays in the volumes 1870-1873 are built by ONE
%  pipeline (ebook-method/for-idee-1870-73/) into a combined working file, which split7.py cuts into seven transcriptions
%  (Aandsvirkelighed and its sequel, Videre om Aandsvirkelighed, are one: Hans, 2026-10-13).
%  Rebuild from the repo root:   sh ebook-method/for-idee-1870-73/rebuild.sh
%
%  STATUS: COMPLETE (2026-10-13): the volume's title page and the whole essay, Bind 4, pp. 289-315, with its notes.
%
%  EDITION.  The only printing.  Text of record: Nasjonalbiblioteket's scan (nb.no) of the volume, ~/bibliotek/Nielsen, Rasmus/1871-1-for-idee-og-virkelighed.pdf (NB item 5afd2e8887e16e47d6f9c4cf0460ebc9) (SCANS.tsv),
%  page images fetched by nbfetch.py, 300 dpi colour, NO text layer; printed page N = PDF page N + 4.  No e-book, no
%  Wikisource text, no reprint, no earlier transcription found (checked 2026-10-13: the repo, the web, the Internet Archive,
%  da.wikisource; HathiTrust and the Google Books API could not be queried that day; Google Books holds at least the 1873
%  volume, nielsen/darwinismen's source -- a second copy, not used).
%
%  ACCURACY (2026-10-13; one collation for all eight texts, 231 pp., ~58,000 words).  The DRAFT is NB's own OCR (ALTO, ABBYY
%  FineReader 12), laid over the page images (combine.py) and turned into LaTeX by draft.py; it was collated word by word
%  with an independent tesseract reading of the same images (collate/check.py, witness 'ocr'): 1556 disagreements.
%  check.py's own rules settled 78.  Candidate rules were first scored BY SCRIPT on a text already read at the image
%  (calibrate.py: Et Synspunkt for Darwinismen, same volume as Verdensanskuelse), then measured on a BLIND PILOT of 152 crops
%  (reader A, 52/52 planted controls right); 17 rules passed (NB's glyph faults: o for ø, D/O/U for Ø, i for l, li for h,
%  l for I, opening » read «, words run together or broken apart; tesseract's: aa read an/ag/å, added accents; the
%  printer's space before a mark, house rule) and settled 772; the rest were read at the image BLIND (562 crops on 94
%  sheets, four readers, 188/188 planted controls right), ~160 decided in the main session from the readers' lines and ~40
%  at 300-dpi zooms.  STRUCTURE read at the image (reader K): every footnote whole (30; collated where they stand, then moved
%  to their marks, then compared word by word with the reading: notecmp.py, 9 differences settled at the image), the
%  headings, the section numerals, the essays' last lines and signatures, the three Greek words, the volume title pages.
%  LETTERSPACING (the print's emphasis, \emph): candidate lines by letter gaps and box width, calibrated on the words the
%  word readers had marked spaced (spaced.py): 361 lines read blind (61/62 roman controls), 25 doubtful extents read at a
%  zoom.  SWEEPS after OPEN 0: words run together that both OCRs share (lower-case + capital), the unknown-word sweep (678
%  words, ~20 suspects, 14 read at the image: 4 OCR faults, 7 misprints kept), --moved, the dropped-word list, the
%  W-shortens list of the blind decisions.
%  SAMPLED PAGE CHECK (pagecheck.py, seed 202610137): one page of each essay or part, pp. 2289, 2400, 3305, 4308, 5033,
%  5208, 6514, 8011 of the working numbering (Bind x 1000 + page), read whole and blind, 1965 words: 0 word errors,
%  1 punctuation (a compound's hyphen lost at a page turn: „Ny-|Holland“) and 1 letterspacing („to be or not to be“)
%  error, both mended.  95 % upper bound for word errors: 1.5 per 1000 words.
%  What this cannot see: a misprint both OCRs normalise into a real word.
%  The ø rule (Hans 2026-10-08): o-forms that are no word are set with ø; „stode“ (= stood) is a word and stays.
%  PRINTED AS IS Bind 4, p. 300: „ligesa lidt“ for ligesaa (zoom z4)
%  Not transcribed: running heads, folios, sheet signatures („For Ide og Virkelighed. 1870. I. 20“), library stamps; the
%  large initial of each essay is set as a plain letter.  Quotation marks as printed (» « in the text, sometimes a closing
%  »; „ “ in some notes).  The notes are numbered ¹), ²) per page, as printed.  Greek as printed (textalpha).
%
%  TOOLS (ebook-method/for-idee-1870-73/): nbalto.py, combine.py, alto.py, ocr.py, draft.py, rules.py, calibrate.py,
%  sheets.py + reads.py + READER-BRIEF.md, spaced.py + spaced.tsv, structsheets.py, notecmp.py, unknown.py, zoom.py,
%  pagecheck.py + PAGECHECK-BRIEF.md, patches.py (PATCHES, PRENOTES, NOTES, POST, MISPRINTS), decisions.tsv, split7.py,
%  rebuild.sh.
% =====================================================================
\documentclass[12pt,a4paper,openany]{book}
\usepackage[utf8]{inputenc}
\usepackage[T1]{fontenc}
\usepackage{libertinus}
\usepackage{libertinust1math}
\usepackage{textalpha}
\usepackage[danish]{babel}
\usepackage[protrusion=true,expansion=true]{microtype}
\usepackage{setspace}
\usepackage[margin=1.3in]{geometry}
\usepackage{fancyhdr}
\usepackage{hyperref}
\hypersetup{pdftitle={Samfundets Naturlove}, pdfauthor={Rasmus Nielsen}, hidelinks}
\pagestyle{fancy}
\fancyhf{}
\fancyhead[LE,RO]{\thepage}
\fancyhead[C]{\textit{R. Nielsen: Samfundets Naturlove (1871)}}
\setstretch{1.3}
\usepackage{marginnote}
\renewcommand{\thefootnote}{\arabic{footnote})}   % the print numbers each page's notes ¹), ²)
\newcommand{\opage}[1]{\marginnote{\footnotesize\textit{[#1]}}}
% the essay's heading: title (and subtitle), centred, a short rule; the print's large initial is not reproduced
\newcommand{\essayhead}[2]{\begin{center}{\large #1}\ifx&#2&\else\\[0.3em]{\normalsize #2}\fi\\[0.5em]\rule{3em}{0.4pt}\end{center}\medskip\noindent}
\newcommand{\signature}[1]{\par\smallskip\hfill #1\hspace*{2em}\par}

\begin{document}

% --- front matter: the title page of Bind 4 (PDF 1 of 1871-1-for-idee-og-virkelighed.pdf), read at the image (reader K, m1);
% short wavy rules are set as plain rules; the library stamps are not transcribed.
\begin{titlepage}
\begin{center}
\vspace*{3em}
{\Huge For Ide og Virkelighed.}\\[1em]
\rule{8em}{0.4pt}\\[1em]
{\large Et Tidsskrift,}\\[0.8em]
{\small udgivet}\\[0.4em]
{\footnotesize af}\\[0.6em]
{\large R. Nielsen,}\\[0.4em]
{\large Bj. Bjørnson og Rud. Schmidt.}\\[1em]
\rule{8em}{0.4pt}\\[0.6em]
{\large Aargang 1871. Første Bind.}\\[0.6em]
\rule{8em}{0.4pt}\\[4em]
Kjøbenhavn.\\[0.4em]
Otto B. Wroblewskys Forlag.\\[1.5em]
{\small G. S. Wibes Bogtrykkeri.}
\end{center}
\end{titlepage}

% ===== TEXT: pp. 289-315 =====
\essayhead{Samfundets Naturlove.}{}
% --- p. 289 ---
\setcounter{footnote}{0}%
\opage{289}Historien har sine Love ligesaa fuldt som Naturen; men
det er vanskeligere at paavise de strenge Love i Historien,
efterdi de ikke umiddelbart gaae op i Nødvendigheden,
men ere Aandslove, mangfoldig betingede
Frihedslove.

Holde vi os til det Almenmenneskelige, saa gaae de
historiske Bevægelser altid ud fra forhaandenværende, til
visse Naturbetingelser bundne Samfundstilstande. At
Samfundet, skjøndt en Frembringelse af Aanden ---
Familieaanden, den patriarkalske Aand o. s. v. --- er
en Naturnødvendighed underlagt, følger ligefrem af, at
Menneskeaanden under alle særlige Fremtrædelser har
sin Grund i Menneskenaturen. Menneskets høiere Natur
er trods Fald og Forbrydelser en aandig Natur; den aandige
Natur er Slægtens Adel. Stamfædrene, Sagnheltene,
de store mythiske Skikkelser, der aabne Historiens Porte
og sætte Begivenhedernes Drivhjul i Gang, ere midt i
den flimrende Fotosfære, den stærke Glands af det Overnaturlige,
hvormed en anende, undrende Fantasi har omgivet
dem, dog i aandsvirkelig Forstand Forbilleder for
Efterkommerne; de ere det netop derved, at det Aandige
% --- p. 290 ---
\setcounter{footnote}{0}%
\opage{290}i deres Daad og Væsen er saa naturligt, og det Naturlige
igjen saa aandigt.

Ifølge Naturretten, det vil sige: Naturens egen faktiske
Ret, gjælder den Grundsætning, at alt det Mulige er tilladt,
og kun det Umulige forbudt. Saaledes forstaaet kan Grundsætningen
ikke anvendes paa den historiske Samfundsudvikling.
I Samfundet er det just det Mulige, der forbydes;
thi at forbyde det Umulige er overflødigt. Her er
i Forhold til Retsfordringerne Meget, der skeer, som
ikke burde skee, og Meget, der burde skee, som ikke
skeer; der er i det menneskelige Samfund, hvad der
ikke er i den blinde Natur, Forskjel og Strid mellem
Ret og Uret.

Ikke destomindre viser Striden overalt hen til en
aandig Naturgrund og forvandler sig, naar Grundkræfterne
røres, til en følelig, haandgribelig Forskjel mellem
det Naturlige og det Unaturlige. Vel er Naturtilstand
ingenlunde Retstilstand; men i afgjørende Forstand
maa dog Retsfordringen støtte sig paa et oprindeligt Krav,
en indre uafviselig Naturtrang. Det Retfærdige er det i
sidste Instans Naturlige, og det Uretfærdige det Unaturlige;
Ret er naturlig, Uret unaturlig; det er disse enfoldige
Grundbestemmelser, hvorpaa det kommer an, naar
Talen er om \emph{Samfundets Naturlove}.

Dersom Sagen nu forholdt sig saaledes, at de Grundfordringer,
der fremgaae af Samfundets Natur, uvilkaarlig
kunde føles og fornemmes af Alle som Fordringer, der
vare naturlige for Individerne og umiddelbart tilfredsstillede
hver Enkelts Fornødenhed: da vilde Overgangen fra
Samfundet i dets Naturstand til en udviklet Retstilstand
foregaae harmonisk, og de historiske Bevægelser være som
% --- p. 291 ---
\setcounter{footnote}{0}%
\opage{291}Tonernes i en Symfoni. Men nu er der, som bekjendt,
en betænkelig Modsætning mellem det, der væsenlig stemmer
med Samfundets Natur, og de umiddelbare, høist
forskjellige Krav hos Individerne. Herfra Brydninger,
Misstemninger, skrigende Dissonanser, der ikke blive opløste.
Istedenfor en naturbegrundet Aandsudvikling fremkommer
nu en skærende Modsætning mellem det Naturlige
og det Konstige, og det Konstige bliver i mange
Maader unaturligt. Til en Opklaring af Samfundets
Naturlove er det nødvendigt at betragte den antydede
Spænding.

Samfundet har ifølge sin Natur mangfoldige baade
stærke og gyldige Krav; men paa Grund af de individuelle
Fordringers endnu større Mangfoldighed komme
disse sociale Naturkrav hvert Øieblik i Strid med hverandre
indbyrdes. Det er en Grundfordring, at der skal
være Orden i Samfundet; men det er ogsaa en Grundfordring,
at Individerne netop i Kraft af denne Orden
maae nyde en vis begrændset Frihed. Der kan til god
Ordens Opretholdelse af de Magthavende anvendes mange
Slags Tvang og Tryk; Individerne finde sig deri og finde
det naturligt for Ordenens Skyld. Alle Velsindede føle,
at der er et Hensyn, der sætter de Enkeltes Privatinteresser
i Skygge, og at det er Hensynet til den offenlige,
statsborgerlige Orden.

Men naar dette Hensyn i længere Tid har gjort sig
ensidigt gjældende, saa bliver Statsmaskineriet mere og
mere konstigt, Anvendelsen af Tryk og Tvang mere og
mere vilkaarlig, og Tilstanden tilsidst aldeles unaturlig.
Samfundet nærmer sig et Vendepunkt, da Millioner
maae lide, for at nogle faa Begunstigede kunne nyde;
% --- p. 292 ---
\setcounter{footnote}{0}%
\opage{292}og naar dette Vendepunkt er indtraadt, naar det ret er
bleven føleligt og vitterligt for Mængden, hvor unaturlig
Tilstanden er: da kan man med samme Sikkerhed forudsige,
at den bestaaende Orden vil blive sprængt, som
man kan forudsige, at der vil skee en Explosion, naar
der gaaer Ild i Krudt. Vel er den Nødvendighed, hvormed
en undertrykt, mishandlet Folkemasse bruser op og
vælter en Throne, uensartet med den Nødvendighed, hvormed
Krudtet exploderer, men Uensartetheden er ikke absolut;
den eneste tilstrækkelige Grund til en Revolution
er en umiddelbar, til uafviselige Livsbetingelser svarende
Nødvendighed, en social Naturnødvendighed.

Frihedens Krav bryder igjennem med Nødvendighed
og anmelder sig derved som en oprindelig Samfundsfordring.
Men naar denne Grundfordring, saa naturlig
den end er, gjøres ensidig gjældende, fremkommer igjen
det Unaturlige. Det er Mængden, der fordrer Friheden
og fordrer den med Lidenskab, men Mængde og Lidenskab
have den medfødte Svaghed tilfælles, at de ikke formaae
at begrændse sig: en fornuftig Selvbegrændsning er
ligesaa fremmed for den vilde Hob som for den opflammede
Lidenskab; den Naturbestemmelse, hvori de begge,
Hoben og Lidenskaben, komme overens, er Tøilesløsheden.
Men for Samfundet er Tøilesløshed det Unaturligste
af Alt; Tøilesløshed er Samfundets Opløsning. Det
er rationel Mekanik. Jo større det Tryk var, der vakte
Reisningen, desto skrækkeligere er Mængdens Rasen; men
jo vildere Mængden raser, desto hurtigere indtræder igjen
Omslaget. Alle Lænker ere brudte, alle Baand løste;
Afgrunden gaber, og Ødelæggelsens Vederstyggelighed vi%
% --- p. 293 ---
\setcounter{footnote}{0}%
\opage{293}ser sig allevegne. Men naar Alt er Kaos, Forvirring og
Rædsel, naar Underverdenen aabner sit Svælg, og Kobbermurene
i Tartarus ere sprængte, hvad Under, at Angst
farer i Mængden, at Skrækken lammer selv de Dumdristige,
at man fra alle Kroge speider forventningsfuld efter
en Jupiter med Lynstraalen. Saa maa der da gives
en styrende Magt, en Magt, der bringer Orden --- Orden
for enhver Pris! Den ordnende Magt kommer udenfra,
ja, ovenfra; men den falder ikke umiddelbart ned fra
Himlen; den har i Tilstanden selv sine Varsler og Forbud.
Despotismen faaer paa slige Vilkaar Naturloven for
sig; Frihedsilden fortærede sig selv: en kraftig Despot begynder Magthaveriet forfra.

Dersom der nu ikke var Andet i Samfundsudviklingen
at opdage end slige, afstikkende Yderligheder, vilde
Historien trods Afvexlingerne i det Udvortes kun være
en Gjentagelse af ensformige Kredsløb, i hvilke den kyndige,
skarpsindige Iagttager næstendels maatte kunne bestemme
Perioderne, som man bestemmer Aarstiderne.
Men ligesom det Aandelige i Mennesket betinges af det
Naturlige, saaledes betinges igjen det Naturlige af det
Aandelige. Hvad Slægten paa eet Samfundstrin kan
finde naturligt, kan den paa et andet finde unaturligt.

Menneskenes naturlige Følelser og naturlige Fordringer
beroe paa en oprindelig Forestilling om, hvad Mennesket væsenlig
er, hvad Værd det har, og hvilket Maal der er sat det i
Verden; af Vexelvirkningen mellem den ideale Forestilling
og den naturlige Selvfølelse i Instinktets og Driftens
Form fremgaaer den urolige Stræben, den aldrig hvilende
Higen mod et fjernt, maaske uopnaaeligt Maal; her er
% --- p. 294 ---
\setcounter{footnote}{0}%
\opage{294}Bevægelsens Aarsag, Grundbetingelsen for alle historiske
Fremskridt.

Det er ikke ligefrem enten Sædernes Forædling eller
Sindets Renselse eller Lidenskabernes Betvingelse, der
danner Grændseskællet mellem de forkristelige og kristelige
Tidsrum i Verdenshistorien. Vel staaer den saakaldte
kristelige Civilisation i mange og væsenlige Træk høit
over den forkristelige, men enkeltvis kunne disse Fortrin
ikke tages til Indtægt. Det er vitterligt, at de kristelige
Dyder have til Følgeskab et Slæng af eiendommelige Laster,
Laster, som den hedenske Oldtid ikke var i Stand til
at fremavle af en saadan Fordærvelsens Dybde, fordi dens
Bevidsthed ikke strakte til; det er vitterligt, at den kirkelige
Tro ikke blot har ægget Vantroen, men endog ved
menneskelig Brøde betinget en Overtro saa mørk og
afskyelig, at Hedenskabets Unaturligheder i Sammenligning
hermed faae et Anstrøg af uvilkaarlig Forblindelse
og undskyldelig Raahed; det er vitterligt, at Aandsforbrydelserne
i alle Retninger voxe med den voxende Bevidsthed
om Aand og desaarsag --- det være sagt, om
man vil, til Ære for den kristelige Civilisation --- aldrig
have rodfæstet sig dybere eller strakt sig videre eller
hævet sig høiere end netop i denne nærværende Tid.

Den naturlige Forestilling om det, at være Menneske,
var i den hedenske Oldtid en ganske anden end
den siden blev ved Kristendommens Indflydelse; det er
Forskjellen i denne naturlige, vel at mærke, \emph{naturlige}
Forestilling om Menneskets Selvværd, der angiver den
historiske Demarkationslinie. Hvor fremtrædende Modsætningen
end er imellem den orientalske Despotisme og
den græske Frihed, saa finder den dog ved nærmere Be%
% --- p. 295 ---
\setcounter{footnote}{0}%
\opage{295}tragtning sin Udligning i et fælles Grundlag. Et Træk
af Historien om Kambyses kan gjøre Sagen anskuelig.
Kongen spørger engang i oprømt Lune sin Ven og Raadgiver
Prexaspes: Hvad sige vel Perserne om mig? ---
Ellers, Herre, svarede Prexaspes, roser man Dig i Alt;
kun mener man, at Du er for hengiven til Vinen. ---
Nu skal Du selv dømme, svarede Kongen, om Perserne
sige sandt, eller om de selv ere fra Forstanden, naar de
sige saadanne Ting. Naar jeg træffer din Søn, som staaer
der i Forsalen midt i Hjertet, saa er det aabenbart, at
Perserne sige usandt, og at jeg er ved min gode Samling.
--- I det Samme spændte han Buen og skjød efter
Drengen, og da han faldt, lod han ham skære op, viste
Faderen, at Pilen sad lige i Hjertet og sagde ret selvtilfreds
til Prexaspes: Altsaa er det jo aabenbart, at det er Perserne,
der ere fra Forstanden, og ikke mig; har Du vel i
den hele Verden seet en bedre Skytte? --- Prexaspes svarede:
En Gud selv kan ikke skyde bedre.

Paa vort nuværende Samfundstrin er der sikkerlig Ingen,
han være forresten nok saa raa og udannet, som
ikke vil finde dénne Handling i høieste Grad oprørende
og unaturlig. Man oprøres ikke blot over Kambyses’s
hensynsløse Grusomhed; man oprøres om muligt endnu
mere over den slaviske Resignation, hvormed Faderen
for at undgaae Sønnens Skæbne, bøier sig under Vilkaarligheden
og smigrer Despotens Almagt. Men Perserne
paa hin Tid med deres Samfundsbevidsthed fandt
ikke en saadan kongelig Handling saa oprørende eller
unaturlig. Tvertimod! De Dannede --- Prexaspes hørte
jo til de Dannede --- maatte i det Hele finde, at Sligt
var i sin Orden. Uden en uindskrænket Magthaver kunde
% --- p. 296 ---
\setcounter{footnote}{0}%
\opage{296}nu engang Samfundet ikke bestaae, og hvad var saa naturligere,
end at den Mægtige kunde handle efter Lune
og hvert Øieblik gjøre, hvad han havde Lyst til.

At dette, hvormange Undtagelser der end kunne
forekomme i det Enkelte, maa have været den herskende
Synsmaade, fremlyser da ogsaa af den mærkelige Kjendsgjerning,
at Kambyses’s Omgivelse yttrede en dyb Medfølelse
for ham selv paa det Sidste. Just som han havde
rustet sig til Kamp imod den falske Smerdis, fik han
ved et Ulykkestilfælde et Saar, der kastede ham paa
Dødsleiet. Han sammenkaldte nu de fornemste Persere,
og for at overbevise dem om, at Oprøreren var en ussel
Løgner, aabenbarede han dem, hvad hidtil var holdt
skjult, at han virkelig havde myrdet sin egen Broder.
Langtfra at angre, beklagede han kun sig selv og begræd
sin sørgelige Skæbne. De Omkringstaaende fandt det
saare naturligt: »da Perserne saae deres Konge græde,
sønderreve de deres Klæder og græd høit med ham».

At slig en Tænkemaade var oprørende for den frie
Græker, er ganske vist; men at det græske Samfund
kunde, om end paa anden Maade og i anden Form, mod
fremragende Individualiteter tillade sig Handlinger, der i
Grunden vare ligesaa oprørende som de, der udmærke Kambyses,
er ogsaa vist. Sagen er, at Grækerne, trods al
deres Sands for Individualitetens frie, harmoniske Udvikling
dog, naar det kom til Stykket, ligesaa lidt som
Perserne havde nogen afgjørende Respekt for Personligheden,
noget umiddelbart gjennemgaaende Indtryk af Menneskets
oprindelige Selvværd. Folkereligionen, det mythiske
Grundlag, paa hvilket den hele, i Skjønhed straalende
Kulturbygning, som vidunderlig smykkede Hellas,
% --- p. 297 ---
\setcounter{footnote}{0}%
\opage{297}blev opført, indeholder en Miskjendelse af Individets evige
Selv, en Fornægtelse af Menneskets personlige Væsen.
Guderne hos Homer udgjøre, som bekjendt, et Samfund
af Udødelige for sig, ophøiet over, og ved en uoverstigelig
Kløft adskilt fra »de dødelige Slægter«. --- »Modige Søn
af Tydeus, hvi spørger Du mig om min Herkomst? Ligesom
Træernes Blade, saaledes er Menneskens Slægter:
Bladene strøes af Blæsten paa Jord, men naar Skovene
grønnes, skyde der friske fra Grenene frem ved Vaarens
Igjenkomst. Saaledes Menneskens Slægter, een fødes, en
anden forsvinder.« Disse Glaukos’s Ord til Diomedes ere
mere end et uvilkaarligt Udbrud af en øieblikkelig Stemning;
de udtrykke en den hele Græcitet gjennemtrængende
Grundbetragtning, en Forudsætning, udover hvilken
den naturlige Selvfølelse aldrig kunde naae. Denne
Miskjendelse af Selvet, denne Bevidsthedens Skranke, var
Ingen af alle Filosoferne i Stand til at hæve; den eneste,
der for Alvor prøvede derpaa, blev dømt til Døden.

Først ved Kristendommen som verdenshistorisk Magt
vinder den personlige Selvfølelse Indgang i Samfundet.
Det er ikke Gudsforholdet i dets Inderlighed og Renhed,
hvortil her sigtes; thi Virkningerne heraf ere ikke ligefrem
udadgaaende og synlige, men indadgaaende, usynlige.
Nei, det er den høiere Forestilling om Menneskets
Bestemmelse i det Hele og den herpaa grundede Følelse
af Menneskeværd og Menneskeret, hvorom Talen er.
Evangeliets Ord om Mennesket som skabt i Guds Billede,
har, naar det lyder i et sandseligt Øre, noget for Selvkjærligheden
høist Indsmigrende; dette Ord har virket
som en Surdeig i Masserne; det har stillet den person%
% --- p. 298 ---
\setcounter{footnote}{0}%
\opage{298}lige Ret i et stedse klarere Lys; det har bragt Selvbetragtningens
Gjenstand nær til Øiet og forstørret Synsbilledet.

Ved at blande det Evige med det Timelige, lade det
Oversandselige træde frem i sandselig Haandgribelighed
og tage de menneskelige Lidenskaber i Tjeneste,
saavel til Begeistring og Handling som til Anger og Lidelse,
formaaede Kirken i Middelalderen at gjøre sig til
øverste stats- og folkestyrende Majestæt. Den havde i
Sammenligning med de lavere Verdensmagter oprindelig
Ret og forstod at gjøre sin Ret gjældende. Til Forstaaelse
af dens Stigen, Kulmination og Dalen er det
nødvendigt at gjøre en skarp Adskillelse mellem Kirken
som religiøs Institution med Salighedsmidlerne i sit Værge
og Kirken som kulturhistorisk Fænomen, som overgribende
Verdensmagt. At Katholicismen i første Henseende har
udfoldet ideale Kræfter og paa sin Vis udviklet stærke
religiøse Karakterer, skal aldrig negtes; men derfor staaer
det ligefuldt fast, at dens høie Selvfølelse som Verdensmagt
var af selvisk Natur; Ilden i dens Blod, Marven i
dens Ben, Drivkraften i dens Hoved og Lemmer var, ogsaa
da, naar den allermest klædte sig i Fromhedens og
Hellighedens Skin, umættelig Havesyge, ubetvingelig Herskesyge,
inkarneret Egoisme. Det er denne Egoisme, der
holder Pavekirken bunden under Naturloven og afstikker
dens Bane. Den steg som Solen, og den maatte stige.
den havde Straaler, som spredte Mørket; den havde et
frelsende, frigjørende Ord; den pegede paa et ophøiet
Maal; den talede i sit Sprog Menneskelighedens Sag, og
dens Ord trængte ind i Sjælene. Men ved selvisk verdslige
Midler at fremme et oververdsligt Maal, at tjene Himlen
ved at herske paa Jorden og tage Ære af Menne%
% --- p. 299 ---
\setcounter{footnote}{0}%
\opage{299}neskene for at udbrede Guds Ære, denne Grundbestræbelse
indeholdt i Spiren en Fordærvelse, der nødvendigt
maatte føre Udviklingen af det første Naturlige over i
det Unaturlige. Tidligt begyndte Bispernes og Ærkebispernes
hierarkiske Anmasselser at falde de nydannede
Stater til Besvær; tidligt begyndte Skandalerne ved det
romerske Hof. Og dog var det Gejstligheden, der tog
sig af de Undertrykkede, støttede Kvinden og frigjorde
Trællen; det var Geistligheden, der i Rettens og Sædelighedens
Navn traadte op imod den brutale Magt og
satte Skranker for barbarisk Vilkaarlighed; det var, med
eet Ord, Geistligheden, der i Næverettens Tid bragte Orden,
og det vildt gjærende Samfund trængte til Orden.

Det er nu let at see, at de samme Betingelser, der
fra Først af medførte Hierarkiets Stigen, maatte i deres
Følger medføre dets Dalen. Dersom det herskesyge, selviske
Væsen blot havde ligget i Individerne, ikke i Kirken
selv, da kunde den indbrydende Fordærvelse, hvor
vidt forgrenet den end til en vis Tid turde være, dog
altid betragtes som noget Tilfældigt og Forbigaaende;
Ingen kunde da med Vished afgjøre, om ikke Pavekirken
maaske lod sig rense og reformere i den Grad, at den
endnu engang fik Styrelsens Tøiler, endnu engang opnaaede
sin gamle middelalderlige Myndighed over Folk og
Fyrster. Men Sagen forholder sig anderledes. Hierarkernes
og Hierarkiets Egoisme er det katholske Kirkevæsens
egen iboende Egoisme. Ved denne, just ved denne
Egoisme, vandt den sit Verdensherredømme, og ved denne,
just ved denne Egoisme tabte den det igjen; Skyggen
er uadskillelig fra Lyset. Det er mærkeligt at see,
hvorledes Styrelsen har holdt det gaaende med denne
% --- p. 300 ---
\setcounter{footnote}{0}%
\opage{300}Magt. Kirken greb det Allerhelligste, opstillede det
som Maal, men brugte det i Virkeligheden som Middel.
At bruge det Allerhelligste som Middel til Opnaaelse af
verdslige Øiemed er Forbrydelse, Aandsforbrydelse efter
den største Maalestok. De, der troe paa Forsynets umiddelbare,
ligefremme Indgriben i Verdensbegivenhedernes
Gang, maae, naar de fæste deres Blik paa de frækkeste
blandt Middelalderens Paver, uvilkaarlig studse og spørge:
»Hvor er det muligt, at Gud kunde tie og saaledes lade
sig spotte; hvor er det muligt, at den Forklarede kunde
sidde saa rolig ved Faderens Høire og see paa, at hans
Billede og hans Værk blev forvrænget, at Herskesygen
ranede Himlens Nøgler og solgte Syndsforladelse for Penge?«
--- Det er Retfærdighed, men en historisk Retfærdighed,
som ikke falder i Øinene. Det Høieste er sat ind;
Menneskene have Frihed til at bruge det; de have ogsaa
Frihed til at misbruge det og tage, hvad der følger paa.
Skulde han, der sidder hos Faderens Høire og følger med
Øiet Verdens Gang, utaalmodigt have knust Pavedømmet
ved et Lynslag, da maatte han jo allerede i sin Tid
have begyndt med at knuse Kajafas og tilintetgjøre Pilatus.

Den historiske Retfærdighed viser sig deri, at Styrelsen,
uden i noget Punkt at opgive sin Ret, overlader
Menneskeslægten den frie Brug af alle de Kræfter og
Midler, der staae til dens Raadighed; det er et Tegn paa
de evige Raadslutningers Uforanderlighed og den guddommelige
Visdoms uendelige Overlegenhed, at Opgaverne
løses, og at Alt dog gaaer naturligt til.

Der er i Historien ligesa lidt Spor af nogen guddommelig   % PRINTED AS IS (Bind 4, p. 300): „ligesa lidt“ for ligesaa (zoom z4)
Forkjærlighed for Papismen som for Muhamedanismen:
Grunden til, at hin kom til at staae høiere
% --- p. 301 ---
\setcounter{footnote}{0}%
\opage{301}end denne, maae søges deri, at den Livskapital, hvormed
den aagrede, var rigere, frugtbarere, ædlere.

For Samfundsudviklingen i det Store gjælder uden
Omsvøb den Grundsætning, at \emph{Ret er Magt}, og Magt
Ret. --- »Og endda skulle vi troe paa et Forsyn?« Ja,
just derfor kunne vi troe paa et faderligt Forsyn; thi et
stærkere Bevis paa Forsynets retfærdige Styrelse kan ikke
gives. Man lægge vel Mærke til, hvad det er, her handles
om. Her handles jo ikke om Samfundets Retfærdighed
mod Individerne i det Enkelte, men kun om den
historiske Ret, hvormed en Samfundsorden til en given
Tid danner sig og hævder sin Bestaaen. Retssætningen
har to Led, af hvilke det ene opklarer det andet. Her
siges ikke, at Magt uden videre er Ret; her siges, at
Magt er Ret, og hvorfor? Fordi \emph{Ret} er \emph{Magt}. Heri
ligger da, at det er Retten, der gjør Magten til Magt;
og dersom denne Bestemmelse gjaldt alene, vilde Fuldkommenhedens
Maal være naaet; det vilde da være Magthaverne
umuligt at misbruge Magten. Men Historien er
i det Timelige, og det Timelige er det Ufuldkomne; saa
gjælder da ogsaa den modsatte Bestemmelse, Ufuldkommenhedens
Udtryk: \emph{Magt er Ret}. Brydningen mellem
de to Grundbestemmelser er paa den mest storartede
Maade blevet anskueliggjort i Pavedømmets Historie.

Her begyndes umiddelbart ud fra den Forudsætning,
at Ret er Magt. Al den menneskelige, naturlige Ret,
hvormed Kirken fra Begyndelsen af saa at sige sad inde,
fik den i Tidernes Løb rig Leilighed til at udfolde. Det
er, som om Historiens Aand pustede til og drev den
frem, at den dog for Himlens Skyld maatte gjøre Alt,
hvad den formaaede, til Menneskeslægtens Opdragelse,
% --- p. 302 ---
\setcounter{footnote}{0}%
\opage{302}til Samfundets Ordning, til Sædelighedens og Kulturens
Fremme. Dette er dens Fart, dens historiske Medbør;
Hierarkiet har Strømmen med sig: fremad, frem! i Vind
og Storm, mellem Brændinger og Skjær gaaer Kirkens
Skib for fulde Seil og kløver Tidens Bølger.

Men da saa Maalet var naaet, da alle de pavelige
Retsbud --- non obstante qvocunque --- vare blevne Magtbud,
hvad saa? Ja, saa begyndte den omvendte Sætning,
Timelighedens Sætning: \emph{Magt er Ret}, at trænge sig
ubændig frem i Begivenhedernes Forgrund. »Det var da
Skade», siger man, »at Kirkeregimentet saaledes skulde
udarte, at Hierarkiet ikke holdt sig i det gamle, ærværdige,
ærefrygtindgydende Spor.« Et tomt Hjertesuk; her er Intet
at sukke over. Hierarkiet udartede, fordi det maatte
udarte, det kunde ikke Andet, det var lovligt undskyldt.
Som en Sten, der kastes i Luften, stiger, men med jævnt
aftagende Hastighed, indtil den naaer et Punkt, da den
saa at sige vender om og begynder at falde med voxende
Hastighed, saaledes er det --- paa Aandens Vei --- gaaet
til med Hierarkiets Stigen og Dalen. Hvad Tyngden er
for Stenen, er Egoismen for Hierarkiet. Først har det
Høiere Overvægt over det Lavere; Ret er Magt: Bevægelsen
gaaer opad; men lidt efter lidt faaer det Lavere
Overvægt over det Høiere; Magt er Ret: Bevægelsen gaaer
nedad. Og ligesom den i Luften svævende Sten har,
naar uvedkommende Paavirkninger fraregnes, en ligesaa
lang Bane at gjennemløbe nedad, som den har gjennemløbet
opad, saaledes har Hierarkiets Fald i Historien beskrevet
en Bane, som ved Krumningernes Afsnit og Længde
giver et omvendt Billede af dets Stigen. Det gik ikke
saaledes til, at Hierarkiet havde en ringere Magt, saa%
% --- p. 303 ---
\setcounter{footnote}{0}%
\opage{303}længe det i det Hele gjorde Magten afhængig af Retten,
og siden fik en umaadelig Magt, da det ret begyndte
at trodse alle Retfærdighedshensyn og gjøre Retten
afhængig af Magten. Tvertimod! den sidste fra Kulminationspunktet
udgaaende Vending er af Historien nøiagtig
betegnet som Magtens begyndende Selvødelæggelse.

For en overfladisk Betragtning tager det sig ud, som
om Pavedømmet blot bukkede under paa Grund af de mange
og forskjellige Slags udvortes Stød og Modstand, det efterhaanden
mødte. Det er en ensidig Opfattelse; den væsenlige
Grund laa i Pavedømmet selv. For at denne
Magt kunde knækkes, maatte den Trin for Trin gjennemgaae
alle Selvnedbrydelsens Mellemstadier. Ikke nok, at
den Begeistringens Ild, der havde opflammet Korsfarerne,
udsluktes; Pavedømmet fik Erstatning: af Troens Æmmer
blussede Overtroen hæftigere op, og Overtroen var
nok saa brugelig som Troen. Ikke nok, at Paven efter
at have fældet den middelalderlige Keiser og derved givet
Middelalderen selv det første Dolkestød, kom i politisk
Afhængighed af Frankrig, thi Afhængigheden blev
kun betalt med Æren: den reelle Profit beholdt Paven
for sig selv: det babyloniske Fangenskab var virkelig indbringende.
 Avignon kunde Pavehoffet udbyde alle Kirkeembeder
og kirkelige Værdigheder offenlig til Auktion;
der blev naturligvis gjort høie Bud ved Auktionerne.
Man skulde ment, at Pavedømmets sidste Time var slaaet,
da man, eftersom der nu var to Pavestole, en i Avignon
og en i Rom, fik paa engang to, siden tre og flere Paver,
der gjensidig fordømte og banlyste hverandre, thi naar
Paven fordømmer Paven, saa har jo Pavedømmet fordømt
sig selv, og Selvmodsigelsen er aabenbar. Men nei!
% --- p. 304 ---
\setcounter{footnote}{0}%
\opage{304}Ogsaa dette Stød lod sig afbøde. For at gjøre en Ende
paa Forargelsen prøvede de frie Koncilier --- i Pisa,
Costnitz og Basel --- ved en Kraftanstrængelse paa at
sætte Paven under et konstitutionelt Formynderskab
og faae en virkelig Reformation i Kirkens Hoved og Lemmer
iværksat. Men dette Skridt, der afgav et slaaende
Beviis for Umuligheden af, at det hierarkisk gjennemførte
System nogensinde skulde kunne reformere sig
selv i evangelisk Aand, faldt netop ud til Pavens Fordel.
Det var da ikke dette eller hint Enkelte, men den
alsidig konsekvente Udfoldning af Magthaveriets Forvrængning,
der gav den kirkelig babyloniske Taarnbygning
sit Grundstød. Just fordi Magten, den blinde Magt,
var bleven Ret, blev Trykket følt i alle Forhold allevegne
og fra alle Sider; Stater, Samfundsklasser, Individer, læge
og lærde, fandt Tilstanden oprørende, utaalelig, unaturlig.
Himlen var overtrukken, Luften kvalm og Horisonten
mørk; Luthers Ord var Lynet i Skyen.

Efterhaanden som det pavelig theokratiske Universalmonarki
adsplittedes og sank, begyndte Statsmagterne i
Europas forskjellige Lande at samle Kræfter og berede
sig til en stor Fremtid. Under Kirkens Formynderskab
kunde Staten ikke komme til fuld Klarhed over sin virkelige
Ret og Betydning. Som pavelige Vasaller vare Fyrsterne
i alle store, afgjørende Anliggender forpligtede til
at lystre Pavens Ordre; en praktisk, naturlig Opfattelse
og Gjennemførelse af virkelige, forhaandenværende Statsinteresser,
var paa disse Vilkaar umulig. Opgaverne vare
verdslige, menneskelige, aldeles naturlige, men Kirken
vidste overalt at underskyde noget Oververdsligt, Overmenneskeligt,
Overnaturligt. Under Paaskud af at adle
% --- p. 305 ---
\setcounter{footnote}{0}%
\opage{305}det Jordiske ved en Tilsætning af det Himmelske, gjorde
den selv det Himmelske til et Jordisk, og dens Indblanding
i Alt blev utaalelig.

Da Reformationens Tordenveir havde renset Luften,
blev Statens Aandedræt friere, ikke blot i de protestantiske,
men ogsaa i de katholske Lande. Staten blev sig,
som det hedder, sin Autonomi (Selvlovgivning) bevidst.
Dette Fremskridt synes ved første Øiekast at være et stort
Tilbageskridt; thi i den klassiske Oldtid --- hos Grækerne
og Romerne --- var jo allerede Staten anerkjendt som det
høieste virkelige Udtryk for den guddommelige Myndighed
paa Jorden: den af Guderne beskyttede Stat var den
i Virkeligheden selv nærværende Guddom. Det maa da
ikke glemmes, at den ud af Middelalderens Gjæringer sig
samlende Stat, den Stat, der tog Magten i Arv efter
Kirken, kom til at staae paa et andet og høiere Grundlag
end den antike; den nyere Stat kan i naturlig og
dog ideal Forstand kaldes den gjenfødte Stat.

Vistnok blev Politiken selv, trods den andægtige
Maske og de Purpurklude af opbyggelige Talemaader,
hvormed den i det sextende og syttende Aarhundrede
saa gjerne smykkede sin Macchiavelisme, ikke
mere kristelig, end den alt havde været i Alexander den
Stores og Keiser Augusts Tidsalder; men den fik i sine
mekaniske Ligevægts- og Bevægelsessystemer nye Masser,
nye Kræfter og nye Betingelser at operere med.
Trods alle Forsyndelser mod det Humane havde den gamle
Kirke dog i Tidernes Løb lidt efter lidt faaet en saa stor
Kapital indsat, om vi saa maae sige, i Humanitetens
Overformynderi, at endog den mest autokratiske Myndig%
% --- p. 306 ---
\setcounter{footnote}{0}%
\opage{306}hed ikke længer ustraffet kunde forholde de Umyndige
deres Ret.

Staten kunde og vilde ikke skyde Kirken fra sig; den
vilde heller ikke paa Oldtidens Vis sætte sig i Kirkens Sted.
Idet den da officielt høitideligt anerkjendte Kirken som
religiøs, relativt selvstændig Institution, der kun i alt
Verdsligt var Statsmyndigheden undergiven, anerkjendte
den med det Samme den Grundopfattelse af Menneskelivet
med de høie ideale Forestillinger, som Kirken i Dogmer,
Ceremonier og Billeder gav Livskraft og sandselig
Nærværelse. Saa blev da Staten nødt til at anerkjende,
at der i hvert Menneskeindivid boer et evigt, personligt
Selv; den blev nødt til at anerkjende, at der foruden de
jordiske og timelige Formaal er sat den hele Slægt et
evigt Maal; at der efter Døden er en Dom; at Alle,
Høie og Lave, Mægtige og Ringe, den kronede Fyrste
som den usleste Tigger, engang skulle møde for den ene
retfærdige Dommer, og Enhver faae efter sine Gjerninger.
Her er et fast, urokkeligt Grundlag. Paa dette Grundlag
kunde nu Magthaveriet med sine politiske Intriger,
Undertrykkelser og Voldshandlinger drive sit Spil saa frit,
det vilde: der var sat Vilkaarligheden en Dæmning.
Gik end Bevægelserne som i en rivende Strøm, en brusende,
svulmende Flod, hvis Bølger her og der overskyllede
de lavere Bredder, druknede Sæden og rev mangen
Blomst og ædel Plante med sig: saa var der dog til
begge Sider en Skranke af Klipper; her var en aandelig
Naturskranke. Vel kunde Masserne endnu paa de mørkeste
Pletter i Samfundet med en vis Ret siges at befinde
sig i en Tilstand af Raahed, Uvidenhed og Overtro,
der ikke gav det gamle Hedenskab Noget efter; men,
% --- p. 307 ---
\setcounter{footnote}{0}%
\opage{307}hvor besynderligt det end lyder, saa var dog Raaheden
selv af en høiere og ædlere Art. Det var en kristelig
Raahed, en kristelig Uvidenhed, en kristelig Overtro. Men
--- siger man --- Raahed, Uvidenhed og Overtro er jo
netop det Ukristelige! Ganske vist. Og dog kalde vi
dette nye Barbari kristeligt, fordi det umuligt kunde
komme frem uden paa kristelige Betingelser: det kunde
kun komme frem i Kristenheden. Det var dette i Kristenheden
Ukristelige, der gjærede i Masserne og satte
Folkene i Uro. Dette i Kristenheden Ukristelige blev en
ny Drivkraft, der gav Bevægelserne i det Store en egen
Retning og paatrykte Følelsen af Taaleligt og Utaaleligt,
\emph{Naturligt og Unaturligt}, sit Særpræg.

Den gigantiske Kraft, hvormed Revolutionen ved det
attende Aarhundredes Slutning brød frem og rev Alting
med sig, vil fra dette Synspunkt vise sig at være begrundet
i en historisk Naturnødvendighed.

Det absolute Monarki, hvori Staten som Stat søgte
at virkeliggjøre sit Ideal, var vel i Stand til at holde de
indbyrdes stridige Elementer i en vis Ligevægt, kue Adsplittelsen
og simplificere Regeringsmaskineriet ved at
samle alle Styrelsens Traade i een Haand; men det var
ikke i Stand til at anerkjende de særlige Institutioners relative
Selvstændighed, ikke i Stand til at forlige Statsmagten
med borgerlige Friheder og almenmenneskelige
Rettigheder. I den enevældige Konges Person blev Staten
selv paa en Maade personlig; men istedenfor at forlige
det Retlige med det Personlige, lod den Afgjørelserne
i høieste Instans beroe paa Naade og Unaade; selv naar
den mente det allerbedst, gav den i Sympathi og Antipathi
Vilkaarligheden Spillerum og satte Personsanseelse
% --- p. 308 ---
\setcounter{footnote}{0}%
\opage{308}høiere end Loven. Staaende paa verdslig Grund udgav
den sig efter Hierarkiets Forbillede for en af Gud selv
umiddelbart indstiftet Orden: det var et kongeligt Theokrati,
der afløste det pavelige. Udgaaende fra Middelalderen,
i hvilken enhver enkelt Stand bar sin egen Fane
saa høit, den kunde, vare nu Stænderne i Samfundet,
Adel, Geistlighed, Borgere og Landalmue --- det var knapt,
at Bondestanden kunde kaldes en Stand --- alle uden
Undtagelse lige for Thronen; men de vare desuagtet langt
fra at være lige for Loven. Det var ikke nok med Individernes
personlige Underkastelse; hver Stand, selv
Adelen, den første, var lige overfor Monarkens allerhøieste
Person en undersaatlig, dybt bukkende, allerunderdanigst
Personlighed. Det er en Selvdrift i dette Personlige, der
først forherliger og derpaa dræber Monarkiet.

Det absolute Monarki var ligesaa lidt en Gjentagelse
af romersk Cæsarisme som af østerlandsk Despotisme ---
Optrin som det med Prexaspes og Kambyses ere en Gang
for alle henviste til det Unaturlige --- den efter Retsprinciper
organiserede Underkastelse af de særskilte, trinvis
opadstigende Personligheder skulde jo netop fyldestgjøre
Samfundets fornuftige Fordringer til en velordnet
Stat. Men i et saadant System af personlige Retshensyn
maatte de særskilte Rettigheder, Begunstigelser, Privilegier
og Forrettigheder dog nødvendig undergrave Retten
som Ret, den almene Ret. Forret for Faa var Uret
mod Mange; og hvori havde saa tilsidst al Forret sit
Støttepunkt, sit egenlige Tyngdepunkt? I Monarkens personlige
Naade: Ret var Naade, og Naade var Ret; deri
var Retten opløst. Det er, lidt dybere seet, ikke de Styrendes
Vilkaarlighed, der misbruger Naadesbevisningerne
% --- p. 309 ---
\setcounter{footnote}{0}%
\opage{309}og derved gjør Naaden vilkaarlig. Sagen forholder sig
omvendt; Naaden er som Retskilde Vilkaarligheden selv:
det er den absolut regerende Naade, der fra Først til
Sidst, fra Øverst til Nederst gjorde Regeringen vilkaarlig.
Det var ingen Tilfældighed, at de enevældige Konger
misbrugte Naaden; nei, det var en aandelig Naturnødvendighed;
er Naaden i sig selv et Misbrug af Retten,
hvor skulde den da kunne bruges uden at misbruges.
Det var, strengt taget, ikke Monarkernes Skyld, at
det absolute Monarki faldt; det faldt --- af kongelig
Naade, efter Naturens Lov: Dødsspiren laa fra Undfangelsen
af i dets Organisme.

Det var den indre, oprørende Splid imellem Retten
som Naade og Retten som Ret, der gjennem utallige
Overgreb, Forbrydelser og Misgreb fremkaldte Katastrofen
i Frankrig. Revolutionen af 1789 blev en Statsomvæltning
i den allerradikaleste Forstand; den førte bogstavelig
til en Sprængning af Staten som Stat. Da Nationalkonventet
afskaffede Kongeværdigheden, og Rædselsherredømmet
tog sin Begyndelse, havde Folkemassen ligefrem
afskaffet Staten og sat sine Kreaturer i Statens Sted; thi
en Robespierre er ingen Statsmand, en Robespierre kan
ikke regere; han kan kun terrorisere. Ved Betragtningen
af et saadant Naturfænomen er det ikke tilstrækkeligt
at dvæle ved Ydersiderne, ved Rædslernes Mangfoldighed,
ved Lidenskabernes Flammehav eller blot ved Lavastrømmenes
forskjellige Udbrudskratere; man maa ned i det
store Svælg, ned i det Indre, forske og spørge: hvor kom
alt dette fra? Hvad var det dog for et Dyb af glødende
Had, hvoraf Ødelæggelsen opstod? Det var --- en kristelig
Revolution, en Revolution, som kun var mulig midt
% --- p. 310 ---
\setcounter{footnote}{0}%
\opage{310}i Kristenheden. --- At Masserne, naar det ydre Tryk bliver
for stort, og Statsmyndighederne under Uroen selv optræde
famlende og splidagtig forraade hverandre, kan i
en hvilkensomhelst Tidsalder ophidses til at nivellere al
Standsforskjel og beruse sig i Indbildninger om Frihed,
Lighed og Broderskab, skal aldrig nægtes. Sligt er muligt
lige saa vel under hedenske som under kristelige
Forudsætninger. Men den Selvfølelsens Styrke, den Instinktets
Sikkerhed, den Tankens Klarhed og Villiens Besluttethed,
hvormed den franske Revolution slog »Menneskerettighedernes«
store Grundlov fast,\footnote[1]{»Det er et karakteristisk Træk, at da den konstituerende Forsamling i Paris (1789) diskuterede Rettighederne, og Mirabeau utaalmodig over, at Diskussionerne antoge et for filosofisk Sving og trak sig i Langdrag, udbrød: N’employez pas le mot de droits, mais dites: Dans l’intérêt de tous, il a été déclaré\dots{} da holdt Forsamlingen med Instinktets Sikkerhed fast paa, at det netop var Rettigheder, altsaa udtrykkelige Principfordringer, der skulde sættes igjennem.» Om Hindringer og Betingelser for det aandelige Liv i Nutiden. S. 34.} lader sig ikke
forklare uden at forudsætte Kristendommen.

Revolutionsmændene rasede mod Kirken og den positive
Religion som overhovedet mod alle historiske Traditioner;
de vilde brænde Forbigangenheden af, rydde
Brandtomterne og begynde forfra med nye Mennesker,
nye Guder og en ny Tidsregning. Men denne revolutionære
Rasen mod Kristendommen er just et talende Bevis
for, hvor inderlig hin mægtige Tidsaand, om end paa
selvisk negativ, dæmonisk Maade havde indsuget og tilegnet
sig Kristendommens første humane Fordringer.

Revolutionens »Frihed og Lighed« er en Forvrængning
af Evangeliets Frihed og Lighed; men ved denne

% --- p. 311 ---
\setcounter{footnote}{0}%
\opage{311}Forvrængning er der ingenlunde gjort et historisk Skridt
til Kristendommens Afskaffelse; der er meget mere stillet
en uafviselig Fordring til dybere Forstaaelse og inderligere
Tilegnelse af Kristendommens Væsen, en Fordring
der lige saa fuldt gjælder Samfundet som den Enkelte.

Modsigelsen mellem »det Nye Testamentes Kristendom«
og  »Kristenheds Kristendom«, en Modsigelse, som
S. Kierkegaard har fremstillet i saa talende, slaaende
Træk, at hans Fremstilling upaatvivlelig engang vil gaae
Verden rundt, denne skærende, aandsfortærende Modsigelse
kræver en dobbelt Løsning, en ethisk-historisk i
det humane Samfund og en, af Menighedslivet betinget,
religiøs Løsning i Samvittighederne. Den sidste skal
ikke afhandles i denne Sammenhæng; her tale vi kun
om den første.

Hvad S. Kierkegaard har kaldt »Kristenheds Kristendom«
er ikke Evangeliet til Salighed, deri har han fuldkommen
Ret; men naar han uden videre betragter den
som et afskyeligt Monstrum, et Uvæsen, der bør afskaffes
jo før jo heller, da har han aabenbart Uret. Kristenheds
Kristendom er trods alle Forvrængninger en ubetvingelig,
historisk Faktor, en Verdensmagt, et Element i den guddommelige
Styrelse; den trygler ikke om Magthavernes
Bifald, den ændser ikke Individernes Vilkaarlighed; den
gaaer, som Stjernerne gaae, Nødvendighedens Vei, dens
Virkninger ere aandige Naturvirkninger. Den følger
med Slægterne; den underlægger sig Tidens Aander; den
leder dem, selv imod deres Villie, stiller dem Opgaverne
og viser dem Fremtidsidealerne som i et magisk
Speil. Hvilke Sider de stillede Opgaver hver Gang frembyde
for Virkeligheden, og hvorledes Synerne i Speilet
% --- p. 312 ---
\setcounter{footnote}{0}%
\opage{312}vexle, beroer efter Tidernes Leilighed paa Evnen, der
virker, og paa Øiet, der seer. Den Slægt, det Samfund,
som med Handlekraft og Sands for Frihed gaaer Selvbegrændsningens
og Maadeholdets Vei, faaer af Kristenhedens
Kristendom aandig Hjælp, mangfoldig Belæring
og en rig Velsignelse. Hvor Kristenhedstanken derimod
ikke renser og forædler Sindelaget, der ægger den og forstærker
de selviske Indbildninger; hvor den ikke kan
virke som en Velsignelse, der virker den som en Forbandelse;
hvor den ikke kan være en Lægedom, bliver
den en Gift.

Den Lov, hvorefter Giftdannelsen foregaaer, lader
sig tydelig paavise. Fra et revolutionært negativt Standpunkt
udslynger Kritiken den Beskyldning mod Kristendommen,
at den, ved at prædike Ydmyghed og blind
Underkastelse under --- Naaden, enerverer Slægten og
stjæler dens bedste Kræfter. Slige Tænkningens Stormløb
laane bestandig deres Retning og Fart fra det, som
angribes. Det urolige Sind vil uden Ophør plage Forstanden,
at den dog endelig maa sætte noget Andet og
Bedre i Stedet for det, som den river ned. Saa skyder
Selvfølelsen iveiret, Masseindividerne stige i Prisen, og
Mængdens Ligeberettigelse til alle Livets Goder lyder
stedse høiere. Det er den foragtede, tilbagetrængte og
fornægtede Kristenhedstanke, der reagerer; som en underjordisk
Ild vedbliver den at hæve den sociale Jordbund
og skrue Livsfordringerne op til Bjerge af Anmasselser.
Hvad der i det gamle hedenske Rom foregik paa Mons
sacer, var kun Børnespil i Sammenligning med det, som
i nærværende Øieblik sees at foregaae paa Montmartre.

Det var sindrigt af Herodot at gjøre Mnemosynes
% --- p. 313 ---
\setcounter{footnote}{0}%
\opage{313}Døttre til sin Histories Muser, thi Mnemosyne betyder
Erindring, og Historie er væsenlig Erindring. Men den
historiske Erindring er af forskjellig Art. For at fremhæve
Natursiden ville vi her see bort fra den Ihukommelse
af Fortidens Bedrifter, som Historieskrivningen
sikkrer, og blot betragte den fysiologisk psykiske Erindring,
den ubevidste halvbevidste Forplantning, hvorved Fædrenes
Dispositioner til Lyder og Dyder nedarves paa
Børnene.

Historien er en Kjæde af sammenhængende Led;
om ogsaa et umiddelbart paafølgende staaer i nok saa
skarp Modsætning til sit foregaaende, saa er der dog i
Bruddet selv et Mærke, som viser, hvor stærk Bindevævet
er. Det Nye fødes, og det Gamle døer, det gaaer til
Grunde. Men, at det Gamle er gaaet til Grunde, betyder,
at det er blevet den Grund, hvoraf det Nye er
opstaaet. Det Efterfølgende har altsaa det Forbigangne
i sin Grund; men hvad det saaledes har i sin Grund,
har det jo paa indvortes Maade i sit Væsen, i sit Indre,
i naturlig Erindring. I det individuelle Menneskeliv er
Ungdommen en Erindring af Barndommen, Manddommen
af Ungdommen, Alderdommen af Manddommen. I Historien
viser sig det Samme. Saaledes er Protestantismen
en Erindring af Katholicismen, Revolutionen en Erindring
af den monarkiske Absolutisme; saaledes er det nittende
Aarhundrede i Tænkemaade, Sæder og Tilstande --- ikke
en Gjentagelse, nei tvertimod --- en Erindring af det
attende. Først naar man har overbevist sig om, at der
i Historien gives en ubevidst, naturnødvendig Erindring,
og at denne Erindring er det indre, usynlige Baand
mellem Slægter og Tidsaldere, da først vil man erkjende,
% --- p. 314 ---
\setcounter{footnote}{0}%
\opage{314}at Menneskehedens Levnetsløb, trods alle tilsyneladende
Standsninger, Gjentagelser og Tilbageskridt, er en ustandselig
Strøm, en fremadskridende Bevægelse mod et sidste
Maal. Det vil da kunne forstaaes, at ikke blot de Dannede,
der have Kundskaber og Overblik, beregne sig en
Fremtid, men at ogsaa Masserne i al deres Raahed og
Uvidenhed have historiske Instinkter.

Vel maa det ved et historisk Jordskælv løbe rundt
for Øinene, naar Samfundsgrunden pludselig revner, og
gamle, længst begravne Magter sees at reise Hovedet
og begynde forfra; man fortvivler om Ret og Fornuft;
man svimler i Kredsløbet og siger med Prædikeren: »Hvad
der har været, det vil være, hvad der er skeet, det vil
skee, og der er intet Nyt under Solen.« Det er igjen
den forslidte Historie med Rigers Reisning og Rigers
Fald, med den Enes Seir og den Andens Nederlag, med
det patriciske Ordensbarbari og det plebejske Frihedsbarbari;
hvor er Fremskridtet? Naar den middelalderlig-romerske
Kirkegud og den gammelromerske Statsgud gjøre fælles Sag,
naar den kanoniske Ret og Kanonretten ere enige om, at
Magt er Ret: hvad Under, at det gaaer ud over de
Afmægtige!

For at kunne holde sig opreist i Jordskælvet og under
den store almindelige Formørkelse skælne Spøgelser fra
Væsener med Kjød og Blod, er det nødvendigt at
støtte sig til Udviklingslovene og see at klare sig den
historiske Erindring. Den pavelige Absolutisme er dømt
til Døden, dømt af Historiens Aand. Dens Tid er forbi;
den er et Gravmonument, et Billede, en Erindringens
Skygge: kommer den igjen, da er det kun som Gjenganger.
Den kongelig-keiserlige Absolutisme er dømt til
% --- p. 315 ---
\setcounter{footnote}{0}%
\opage{315}Døden, dømt af Historien selv; den kan reise sig op,
men ikke skride frem: den tynges af Forgangenhedens
Lænker. Vender Gjenfærdet tilbage fra sit Skyggerige,
da er det som Knokkelmand med et Slagsværd til Le
og et Dødninghoved under Hjelmen: Ridderen af den sorte
Ørn kommer ikke for at saae, men for at høste. Den
demokratisk-revolutionære Absolutisme er dømt til Døden,
og hører alt til de Dødes Rige. »Det røde Spøgelse«
paa Barrikaderne har ikke Blod i Aarerne, ikke Veir i
Lungerne, det har ikke Nærværelsens Liv: seet i Mørket
er det et rædsomt, bloddryppende Fantom, men det er
kun et Fantom.

For at bestemme Himmellegemernes Baner, maa
man beregne Perturbationerne. De virkelige Baner ere
jo vistnok meget afvigende fra dem, der svare til de
abstrakt-almene Love; men disse Afvigelser ere saa langt
fra at vække Tvivl om Grundlovenes Gyldighed, at de i
Virkeligheden stadfæste deres Opfyldelse. De Perturbationer,
som Fænomenerne af de store Samfundslove og
de historiske Baner ere underkastede, lade sig vel ikke
saaledes mathematisk beregne; men derfor er det ligefuldt
sandt, at de alle gaae op i Lovenes Opfyldelse. Er her
mange ubekjendte Størrelser, saa er her ogsaa mange
Ligninger. Endeligningen er oprindeligt bestemt ved
den i sin Grund konstante Menneskenatur og den under
alle Forandringer uforanderlige Kristustanke.
\signature{R. \emph{Nielsen}.}

\end{document}
